% ===================================================================
%  કોષ્ટક નં. ૩ — બાળપણ અને યુવાની.
%
%  Traduction, faite en 2026, en gujarati d'aujourd'hui, sur l'ido de
%  texto/io. Regles de la colonne : voir 10-kostak-01.tex. Deux
%  scenes, deux numerotations : t03-c1-* le bapteme au village,
%  t03-c2-* la fete publique. Ecrit parigi.py ouvert : vingt paires
%  annoncees pour ce seul tableau, dont les quatre du bloc
%  t03-c1-23-1.
%
%  LE « julio » DU BLOC t03-c2-01-1. Le fac-simile ido compose
%  « (julio od \VUgras{agosto}) » en laissant le seul juillet hors du
%  gras. gravuri/korekti.json le repare pour chaque colonne, et le
%  triple gujarati — « (જુલાઈ અથવા » — a ete pose au cablage, AVANT
%  que ce tableau soit ecrit.
%
%  CORRECTION, ET ELLE VAUT POUR DEUX MESSAGES DE COMMIT. Le cablage
%  de cette colonne annoncait « corrections declarees : 61 a 62 », et
%  celui de la colonne haoussa « 60 a 61 ». Les deux etaient faux
%  d'un tableau : html.py ne compte pas les triples DECLARES dans
%  korekti.json, il compte ceux qu'il a effectivement APPLIQUES, et
%  un triple pose pour le tableau 3 ne s'applique qu'une fois le
%  tableau 3 ecrit. Le compte passe donc a 62 ICI, et il etait passe
%  a 61 au tableau 3 haoussa. Poser la correction d'avance reste la
%  bonne facon de faire ; c'est l'annoncer comme deja comptee qui
%  etait une erreur.
%
%  LE GUJARATI A SIX SAISONS, ET LA PLANCHE N'EN CONNAIT QUE QUATRE.
%  Le releve en a maintenant trois comptes differents : quatre en
%  Europe, cinq au Sahel — rani, bazara, damina, kaka, hunturu —, et
%  six ici : વસંત le printemps, ગ્રીષ્મ l'ete, વર્ષા la mousson, શરદ
%  l'automne, હેમંત l'avant-hiver, શિશિર l'hiver. La saison des
%  pluies est celle qu'aucune des deux autres n'a comme telle, et
%  c'est celle qui fait toute la difference : la ou l'Europe partage
%  l'annee par la temperature, l'Inde la partage par l'eau. Le
%  printemps de ce tableau tombe juste sur વસંત, mars-avril-mai.
%
%  ET LES SEPT JOURS SONT LES SEPT PLANETES, EXACTEMENT COMME EN
%  EUROPE. રવિવાર le jour du Soleil, સોમવાર celui de la Lune, મંગળવાર
%  celui de Mars, બુધવાર de Mercure, ગુરુવાર de Jupiter, શુક્રવાર de
%  Venus, શનિવાર de Saturne. Ce n'est pas une ressemblance : c'est la
%  MEME semaine, sortie de la meme astrologie hellenistique, arrivee
%  en Inde par les Grecs de Bactriane et a Rome par l'Egypte. Le
%  releve a donc trois reponses a la meme question — le haoussa a
%  emprunte ses sept jours a l'arabe, le persan en a numerote six et
%  n'a laisse a l'arabe que le vendredi, le gujarati a garde les
%  planetes — et la troisieme est la seule ou le francais et la
%  colonne disent litteralement la meme chose.
%
%  QUATRE FAUTES D'ORDRE, ET parigi.py LES AVAIT TOUTES LES QUATRE.
%  Ce n'est pas le tableau 1 qui se repete : la liste avait ete lue
%  cette fois, et les quatre blocs y sont — t03-c1-04-1, c1-09-1,
%  c1-10-1, c1-14-1. Ce qui manquait n'etait pas l'avertissement,
%  c'etait le TOUR DE PHRASE. Savoir que « (7) doit preceder (5) »
%  ne sert a rien tant qu'on n'a pas la construction qui le permet
%  dans une langue ou le determinant precede le determine.
%
%  LE REMEDE EST LA RELATIVE, ET IL VAUT POUR TOUTE LANGUE A TETE
%  FINALE. Le gujarati a une relative postposee en « જે... », heritee
%  du sanskrit, qui pose la TETE d'abord et rejette derriere tout ce
%  qui la qualifie — c'est-a-dire exactement ce que l'ordre de la
%  planche demande. « અબાબીલ ... શિખરની આસપાસ ઊડે છે, જે ઘંટમંદિરનું
%  છે અને જે દેવળને મુગટ પહેરાવે છે » rend 4, 7, 5, 6 sans forcer la
%  langue d'un pouce. Les quatre blocs se sont corriges par le meme
%  moyen. parigi.py dit QUELLES paires ; la relative dit COMMENT.
%
%  L'EGLISE EST PORTUGAISE, ET C'EST LA QUATRIEME REPONSE DU RELEVE.
%  Le persan nommait le bapteme en arabe et l'eglise en syriaque ; le
%  haoussa nommait tout en anglais, coci, baftisma, fada ; le
%  gujarati dit પાદરી le pretre, du portugais PADRE, et બાપ્તિસ્મા —
%  parce que le christianisme est entre en Gujarat par Diu et Daman,
%  et non par Londres. C'est la meme couche que « મેજ » la table du
%  tableau 1, et elle traverse toute la colonne.
%
%  SECOND REFUS DE LA COLONNE, ET QUATRIEME DU RELEVE SUR CE MEME
%  RENVOI : LE BAL PUBLIC (42). Le javanais avait « joged » et ne l'a
%  pas ecrit ; le haoussa avait « bori » et ne l'a pas ecrit ; le
%  gujarati a « ગરબા », la ronde de Navratri, neuf nuits, un cercle
%  autour d'une lampe, une langue entiere de pas et de battements de
%  mains. La planche grave un bal de sous-prefecture avec estrade et
%  lampions. On ecrit « જાહેર નૃત્ય ». TROIS COLONNES, TROIS DANSES
%  NATIONALES, TROIS REFUS AU MEME NUMERO.
% ===================================================================

%%K t03-apar-1 apar
\VUsaut{3.0mm}
\VUtitre{15.0pt}{60}{કોષ્ટક નં. ૩}
\VUsaut{1.6mm}
\VUtitre{11.4pt}{0}{બાળપણ અને યુવાની.}
\VUsaut{3.4mm}

%%K t03-apar-2 apar
(ત્રીજું અને ચોથું કોષ્ટક એવી રીતે જોડાયેલાં છે કે તે ચાર
\VUgras{કૌટુંબિક દૃશ્યોમાં} \VUgras{હવામાન,} \VUgras{ઉંમર,}
\VUgras{કુટુંબ,} \VUgras{પોશાક} અને \VUgras{સ્થિતિ} વિશેના પાયાના
ખ્યાલો રજૂ કરે છે.)

\VUsaut{4.0mm}
%%K t03-tit-1 sub
\VUcentre{12.0pt}{0}{\textit{પ્રથમ દૃશ્ય.}}
\VUsaut{3.0mm}

%%K t03-tit-2 sub
\VUpk{10.2pt}{40}{ગામમાં બાપ્તિસ્માની વિધિ.}
\VUsaut{2.4mm}

%%K t03-tit-3 sub
\VUpk{10.2pt}{40}{વસંત.}
\VUsaut{3.0mm}

%%K t03-c1-01-1 p
1. --- અમે \VUgras{વસંતમાં} છીએ, \VUgras{માર્ચ,} \VUgras{એપ્રિલ} કે
\VUgras{મે} મહિનામાં, \VUgras{અઠવાડિયાના} દિવસોમાં,
\VUgras{સોમવારે,} \VUgras{મંગળવારે} કે \VUgras{બુધવારે.} સવારના
\VUgras{દસ વાગ્યા} છે.

%%K t03-c1-02-1 p
2. --- \VUgras{હવામાન} સુંદર છે, \VUgras{હવા} સ્વચ્છ છે. સૂરજ
ચમકે છે. કેટલાંક \VUgras{નાનાં વાદળ} \textsuperscript{(2)} ભૂરા
\VUgras{આકાશમાં} \textsuperscript{(1)} ચઢે છે.

%%K t03-c1-03-1 p
3. --- \VUgras{ઝાડને} હજી ભાગ્યે જ \VUgras{પાંદડાં} છે. પાતળાં
\VUgras{પોપ્લર} \textsuperscript{(3)} અને ઊંચું \VUgras{ચિનાર}
\textsuperscript{(17)} અત્યાર સુધી માત્ર કળીઓ જ ધરાવે છે.

%%K t03-c1-04-1 p
4. --- \VUgras{અબાબીલ} \textsuperscript{(4)} પાછાં આવે છે અને
આનંદથી કલરવ કરતાં \VUgras{શિખરની} \textsuperscript{(7)} આસપાસ ઊડે છે,
જે \VUgras{ઘંટમંદિરનું} \textsuperscript{(5)} છે અને જે ગામના
\VUgras{દેવળને} \textsuperscript{(6)} મુગટ પહેરાવે છે.

%%K t03-tit-4 sub
\VUsaut{3.4mm}
\VUpk{10.2pt}{40}{દેવળ.}
\VUsaut{3.0mm}

%%K t03-c1-05-1 p
5. --- આ દેવળ પોતાના મોટા \VUgras{દરવાજા} \textsuperscript{(13)} અને
મોટા \VUgras{ઘડિયાળ} \textsuperscript{(8)} સાથે ઘણું સાદું છે. નાનો
\VUgras{કાંટો} \textsuperscript{(9)} X ના આંકડા પર છે, તે
\VUgras{કલાક} બતાવે છે. \VUgras{મોટો} \textsuperscript{(10)} XII પર છે
(બપોર) ; તે \VUgras{મિનિટ} બતાવે છે.

%%K t03-c1-06-1 p
6. --- ઘંટમંદિરમાં \VUgras{ઘંટ} \textsuperscript{(11)} આનંદથી વાગે છે.
છાપરાની ટોચ પર પથ્થરનો નાનો \VUgras{ક્રૂસ} \textsuperscript{(12)} ઊભો
છે.

%%K t03-c1-07-1 p
7. --- દેવળને માત્ર મોટો \VUgras{દરવાજો} \textsuperscript{(13)} જ નથી,
તેને બાજુમાં નાનું બારણું પણ છે. આ નાના બારણામાંથી પોતાનો
\VUgras{ઝભ્ભો} પહેરેલા \VUgras{પાદરી} \textsuperscript{(15)} સાહેબ
\VUgras{વસ્ત્રખંડમાંથી} \textsuperscript{(14)} બહાર નીકળીને
\VUgras{પાદરીના ઘર} \textsuperscript{(19)} તરફ જાય છે.

%%K t03-c1-08-1 p
8. --- તેની પાસે, આ રહી \VUgras{દૂધ વેચનારી}
\textsuperscript{(24)} પોતાના \VUgras{ગધેડા} \textsuperscript{(23)} સાથે.
તે શહેરમાંથી પાછી આવે છે, જ્યાં તે બાળકો માટે સારું દૂધ લઈ ગઈ હતી.

%%K t03-tit-5 sub
\VUsaut{4.0mm}
\VUpk{10.2pt}{40}{ફળનો બગીચો.}
\VUsaut{3.4mm}

%%K t03-c1-09-1 p
9. --- ગધેડાભાઈ આ સવારની દોડથી તદ્દન સંતુષ્ટ છે. તે
\VUgras{ફૂલોની} \textsuperscript{(20)} સુગંધથી ભરેલી હવા મજાથી સૂંઘે
છે, જે ફૂલો પડોશના \VUgras{ફળના બગીચાનાં}
\VUgras{ફળઝાડને} \textsuperscript{(18)} ઢાંકે છે.
આ \VUgras{આલૂનાં ઝાડ} \textsuperscript{(18)} અને \VUgras{જરદાળુનાં ઝાડ}
\textsuperscript{(19)} તદ્દન ગુલાબી અને સફેદ છે, અને તે સારા પાકની આશા
આપે છે.

%%K t03-c1-10-1 p
10. --- એ દરમિયાન નાની \VUgras{મધમાખીઓ}
\textsuperscript{(22)} \VUgras{મધપૂડામાંથી} \textsuperscript{(21)}
નીકળીને ગુંજતી ગુંજતી ત્યાં પોતાનું મીઠું \VUgras{મધ} ભેગું કરવા
જાય છે.

%%K t03-c1-11-1 p
11. --- પણ શું થાય છે ? આ રહ્યા ગામનાં છોકરાં દોડી આવે છે અને ડરી ગયેલાં
\VUgras{હંસને} \textsuperscript{(25)} ભગાડે છે. નોટરીના દીકરાના
\VUgras{બાપ્તિસ્મા} પછી દેવળમાંથી નીકળતું સરઘસ છે.

%%K t03-c1-12-1 p
12. --- નાનકડો યાકોબુસ, પોતાના \VUgras{પારણામાં} \textsuperscript{(26)},
ઘંટના આનંદી નાદથી જાગી જાય છે, અને તેની બહેન માગ્દાલેના, જે પોતે પણ
બાપ્તિસ્માનું સરઘસ જોવા માગે છે, પોતાની માને વિનંતી કરે છે કે તે
આવીને ઘરના ઉંબરે તેના વાળ ઓળે અને તેને કપડાં પહેરાવે.

%%K t03-c1-13-1 p
13. --- મા સંમત થાય છે. તે પોતાનો \VUgras{રૂમાલ}
\textsuperscript{(28)}, પોતાનું \VUgras{કાપડું} \textsuperscript{(29)}
અને પોતાનું \VUgras{એપ્રન} \textsuperscript{(30)} પહેરે છે, અને બારણા
આગળ નાની ખુરશી પર બેસવા આવે છે. માગ્દાલેનાએ તો માત્ર પોતાનું
\VUgras{ઘાઘરું} \textsuperscript{(31)} અને પોતાનું \VUgras{કાંચળી}
\textsuperscript{(32)} જ પહેર્યાં છે.

%%K t03-c1-14-1 p
14. --- તે કેટલી ખુશ છે ! તે પોતાનાં પ્રિય ફૂલ ભૂલી જાય છે, પોતાનાં
\VUgras{કાર્નેશન} \textsuperscript{(34)}, જેમના પર \VUgras{પતંગિયાં}
\textsuperscript{(35)} બેસે છે, પોતાનાં \VUgras{ટ્યૂલિપ}
\textsuperscript{(36)}, પોતાનાં \VUgras{મેનાં ફૂલ}
\textsuperscript{(37)}, \VUgras{કૂંડાંમાં} \textsuperscript{(38)},
\VUgras{ભીંત} \textsuperscript{(33)} પર, અને પોતાનાં \VUgras{ગુલાબ}
\textsuperscript{(39)}, જે એટલી સુંદર વાડ બનાવે છે. તે સરઘસની
સ્ત્રીઓનાં કીમતી કપડાં જુએ છે.

%%K t03-c1-15-1 p
15. --- ગામનાં છોકરાં કપડાં તરફ બહુ ધ્યાન આપતાં નથી. તેઓ
\VUgras{ચોકલેટ} \textsuperscript{(55)} કે \VUgras{સિક્કા}
\textsuperscript{(52)} મેળવવા માટે ધસી જાય છે અને એકબીજાને ધક્કો મારે
છે. તેમાંનું એક રડે છે \textsuperscript{(40)}.

%%K t03-c1-16-1 p
16. --- બીજાએ \VUgras{કૂતરાના} \textsuperscript{(42)} પગ પર પગ મૂક્યો,
જે રડતો રડતો ભાગે છે અને \VUgras{મરઘીને} \textsuperscript{(43)} અને
તેનાં \VUgras{બચ્ચાંને} \textsuperscript{(44)} ભગાડે છે.

%%K t03-tit-6 sub
\VUsaut{5.0mm}
\VUpk{10.2pt}{40}{બાપ્તિસ્માનું સરઘસ.}
\VUsaut{4.0mm}

%%K t03-c1-17-1 p
17. --- સરઘસની આગળ \VUgras{ધાવ} \textsuperscript{(45)} ચાલે છે. તેને
લાંબી રિબનવાળી સુંદર \VUgras{ટોપી} \textsuperscript{(46)} છે. તે
\VUgras{નવજાત બાળકને} \textsuperscript{(47)} ઊંચકે છે, જેને તદ્દન
\VUgras{ભરતકામથી} \textsuperscript{(48)} વીંટ્યું છે.

%%K t03-c1-18-1 p
18. --- ધાવ પછી \VUgras{ધર્મપિતા} \textsuperscript{(49)} અને
\VUgras{ધર્મમાતા} \textsuperscript{(53)} આવે છે. તેઓ ચોકલેટ અને સિક્કા
વહેંચે છે. ધર્મપિતાએ \VUgras{મોજાં} \textsuperscript{(51)} અને
\VUgras{ઊંચી ટોપી} \textsuperscript{(50)} પહેર્યાં છે. ધર્મમાતા હાથમાં
સુંદર \VUgras{ફૂલગુચ્છ} \textsuperscript{(54)} પકડે છે.

%%K t03-c1-19-1 p
19. --- તેમની પાછળ, અમે બાળકના \VUgras{કાકા} \textsuperscript{(56)}
અને \VUgras{ફોઈ} \textsuperscript{(58)} ને જોઈએ છીએ. ફોઈએ સુઘડ
\VUgras{ટોપી} \textsuperscript{(59)} પહેરી છે, કાકાએ કાળો
\VUgras{લાંબો કોટ} \textsuperscript{(57)} અને સફેદ બંડી. પછી આ રહ્યા
\VUgras{કાકાનો દીકરો} \textsuperscript{(60)} અને \VUgras{કાકાની દીકરી}
\textsuperscript{(61)}. તે હાથ પકડીને નવજાત બાળકની \VUgras{બહેન}
\textsuperscript{(62)} ને, નાનકડી બેર્તાને, દોરે છે.

%%K t03-c1-20-1 p
20. --- બેર્તાનો બીજો \VUgras{ભાઈ} \textsuperscript{(63)} પાછળ ચાલે છે.
તે એક \VUgras{સગી} \textsuperscript{(64)} નો હાથ પકડે છે. કુટુંબનાં
\VUgras{મિત્રો} \textsuperscript{(65)} પછી આવે છે.

%%K t03-tit-7 sub
\VUsaut{4.6mm}
\VUpk{10.2pt}{40}{જોનારાં.}
\VUsaut{3.4mm}

%%K t03-c1-21-1 p
21. --- સરઘસની પાસે, આ રહ્યો \VUgras{ભિખારી} \textsuperscript{(66)}
પોતાની \VUgras{ઘોડી} \textsuperscript{(67)} પર ટેકવીને. તે ભીખ માગે છે.

%%K t03-c1-22-1 p
22. --- બીજી બાજુએ ઘણો \VUgras{વિનયી} \textsuperscript{(68)} નાનો
છોકરો છે. તે પોતે જેમને ઓળખે છે તેમને નમસ્કાર કરે છે અને
« \VUgras{સુપ્રભાત} » કહે છે.

%%K t03-c1-23-1 p
23. --- ગામલોકો બાપ્તિસ્માનું સરઘસ જોવા પોતાનાં ઘરમાંથી બહાર નીકળ્યાં
છે. અમે \VUgras{ખેડૂતને} \textsuperscript{(69)} પોતાની \VUgras{સુતરાઉ
ટોપી} સાથે અને તેની બાજુમાં \VUgras{ખેડૂતાણીને} \textsuperscript{(70)}
જોઈએ છીએ. પછી \VUgras{ડોસી} \textsuperscript{(71)} પોતાની
\VUgras{લાકડી} \textsuperscript{(72)} પર ટેકવીને. છેલ્લે \VUgras{ઘંટીવાળો}
\textsuperscript{(73)} પોતાની મોટી \VUgras{નમદાની ટોપી}
\textsuperscript{(74)} સાથે અને \VUgras{લાકડાનાં ચંપલ}
\textsuperscript{(75)} પહેરેલી જુવાન ખેડૂતાણી, જે પોતાના બાળકને ચાલતાં
શીખવે છે.

%%K t03-c1-24-1 p
24. --- આ દૃશ્ય ઘણું મજાનું છે.

\VUsaut{4.0mm}
%%K t03-tit-8 sub
\VUcentre{12.0pt}{0}{\textit{બીજું દૃશ્ય.}}
\VUsaut{3.0mm}

%%K t03-tit-9 sub
\VUpk{10.2pt}{40}{જાહેર ઉત્સવ.}
\VUsaut{2.4mm}

%%K t03-tit-10 sub
\VUpk{10.2pt}{40}{નૌકારમતો અને હરીફાઈઓ.}
\VUsaut{3.0mm}

%%K t03-c2-01-1 p
1. --- \VUgras{ઉનાળો} આવી ગયો છે. \VUgras{જૂન} મહિનાનો (જુલાઈ અથવા
\VUgras{ઓગસ્ટ}) સખત તડકો જુવાનોને નહાવા અને હોડી ચલાવવા પ્રેરે છે.

%%K t03-c2-02-1 p
2. --- નદીના જમણા કાંઠે, હોડી ચલાવનારા નૌકારમતો અને હરીફાઈઓમાં ભાગ
લેવા માટે કપડાં ઉતારે છે.

%%K t03-c2-03-1 p
3. --- તેમાંનો એક પોતાનાં \VUgras{ચંપલ} \textsuperscript{(1)} ઉતારે છે ;
તે તેમની દોરી છોડે છે. તેના બે સાથી તૈયાર થઈ ગયા છે. હવે તેમની પાસે
માત્ર પોતાનું \VUgras{ગૂંથેલું બંડી} \textsuperscript{(2)} અને
\VUgras{નહાવાની ચડ્ડી} \textsuperscript{(3)} જ છે. તેઓ પોતાના મિત્રોની
રાહ જોતાં વાતો કરે છે. તેઓ રજા વિશે અને સામા કાંઠે થતા ઉત્સવ વિશે વાત
કરે છે.

%%K t03-c2-04-1 p
4. --- જમીન પર, તેમની પાસે, તેમના સાથીઓનાં \VUgras{કપડાં} પડ્યાં છે :
\VUgras{બૂટની} \textsuperscript{(4)} જોડ, \VUgras{ચંપલ}
\textsuperscript{(5)}, \VUgras{મોજાં} \textsuperscript{(6)},
\VUgras{નમદાની} \textsuperscript{(7)} અને \VUgras{ઘાસની}
\textsuperscript{(8)} ટોપીઓ અને \VUgras{ટાઈ} \textsuperscript{(9)}.

%%K t03-c2-05-1 p
5. --- જુવાનોમાંના એકે પોતાનાં \VUgras{કપડાં} \textsuperscript{(10)}
કાળજીથી વાળ્યાં છે. તે તેમને એક ટુવાલ પર મૂકે છે, જેમાં તેનો પડોશી
થોડી વારમાં પોતાનાં બે \VUgras{કપડાં} બાંધશે.

%%K t03-c2-06-1 p
6. --- છઠ્ઠો છોકરો મોડો પડે છે. તેના માથે હજી પોતાની \VUgras{ટોપી}
\textsuperscript{(11)} છે. તેણે પોતાનું \VUgras{પહેરણ}
\textsuperscript{(12)}, પોતાનો \VUgras{કોલર} \textsuperscript{(13)} કે
પોતાનાં \VUgras{કફ} \textsuperscript{(14)} ઉતાર્યાં નથી. તે હાથમાં
પોતાની \VUgras{બંડી} \textsuperscript{(17)} પકડે છે અને હજી પોતાનું
\VUgras{પાટલૂન} \textsuperscript{(16)} અને પોતાના \VUgras{ગાળિયા}
\textsuperscript{(15)} પહેર્યાં છે. તે આવતી હરીફાઈ માટે તૈયાર નહીં થાય
એ નક્કી.

%%K t03-c2-07-1 p
7. --- જુવાનોની પાસે, એક નાની કેડી, જેની બંને બાજુએ ઘણાં
\VUgras{કાંટાળાં ઝાંખરાં} \textsuperscript{(18)} છે, નદીના કાંઠે લઈ
જાય છે, જ્યાં યાકોબુસ કાકા \VUgras{બરુ} \textsuperscript{(19)} વચ્ચે
\VUgras{દારૂખાનું} \textsuperscript{(20)} તૈયાર કરે છે, જે સાંજે
ફોડવામાં આવશે.

%%K t03-tit-11 sub
\VUsaut{4.4mm}
\VUpk{10.2pt}{40}{હરીફાઈ.}
\VUsaut{3.4mm}

%%K t03-c2-08-1 p
8. --- નદી પર, કેટલીક સ્ત્રીઓ, પોતાની છત્રીથી પોતાને છાંયડો કરતી,
\VUgras{હોડીમાં} \textsuperscript{(21)} ફરે છે. તેઓ ઘણી રસપ્રદ હરીફાઈ
જુએ છે. દરેક \VUgras{નાવડીમાં} \textsuperscript{(22)} ચાર
\VUgras{હલેસાં મારનારા} \textsuperscript{(24)} પોતાની બધી તાકાતથી હલેસાં
મારે છે, જ્યારે \VUgras{સુકાની} \textsuperscript{(26)}
\VUgras{સુકાન} \textsuperscript{(25)} પકડે છે અને નાજુક હોડીને દોરે છે.
\VUgras{હલેસાં} \textsuperscript{(23)} તાલમાં પાણી પર પડે છે અને હલકી
હોડીઓ ચક્કર આવે એવી ઝડપે ચાલે છે.

%%K t03-c2-09-1 p
9. --- કેટલાક જુવાનો, પુલના થાંભલા પાસે, \VUgras{લપસણા થાંભલાના}
\textsuperscript{(28)} ઇનામ માટે હરીફાઈ કરે છે. તેમાંનો એક છેડે બાંધેલી
નાની \VUgras{ધજા} \textsuperscript{(29)} સુધી પહોંચવા માટે લવચીક અને
લપસણા થાંભલા પર સાવધાનીથી ચાલે છે. તેનો કમનસીબ \VUgras{હરીફ}
\textsuperscript{(30)} પાણીમાં પડ્યો છે. તે કિનારે પહોંચવા તરે છે.

%%K t03-c2-10-1 p
10. --- મોટી ઉંમરના જુવાનો \VUgras{ધક્કાધક્કીની હરીફાઈ}
\textsuperscript{(32)} \VUgras{ધક્કા} \textsuperscript{(33)} પાસે કરે છે.
આ બધી રમતો મોટી \VUgras{જનતા} \textsuperscript{(31)} જુએ છે,
\VUgras{પુલની} \textsuperscript{(27)} \VUgras{પાળી} પર ટેકવીને અને
કાંઠા પર ભેગી થઈને. તેમ છતાં કેટલાક જોનારા \VUgras{ઇનામી થાંભલાની}
\textsuperscript{(45)} રમત જોવા માટે કે ઇનામ \VUgras{વહેંચવા} આવતા
\VUgras{નગરપતિના સરઘસને} વધાવવા માટે પાછા ફરે છે.

%%K t03-c2-11-1 p
11. --- પહેલાં, આ રહ્યા કેટલાક \VUgras{છોકરાઓ} \textsuperscript{(35)}
\VUgras{વાજિંત્રવાળાઓની} \textsuperscript{(36)} આગળ ચાલતા ; પછી
\VUgras{નગરપતિ} \textsuperscript{(37)}, તેમના મદદનીશો અને નાના શહેરની
\VUgras{નગરસભા} આવે છે. છેલ્લે \VUgras{અગ્નિશામકો}
\textsuperscript{(38)} ચાલે છે.

%%K t03-tit-12 sub
\VUsaut{5.0mm}
\VUpk{10.2pt}{40}{મેળો.}
\VUsaut{3.6mm}

%%K t03-c2-12-1 p
12. --- \VUgras{મેળાના મેદાન} \textsuperscript{(39)} પર ઘણા લોકો છે.
લોકો જુદા જુદા \VUgras{માંડવા} જોવા આવે છે : \VUgras{લાકડાના ઘોડા}
\textsuperscript{(40)}, \VUgras{સરકસ} \textsuperscript{(41)}, જ્યાં ખેલ
શરૂ થઈ ગયો છે ; \VUgras{જાહેર નૃત્ય} \textsuperscript{(42)}, જ્યાં
શહેરના જુવાનો અને જુવાનડીઓ \VUgras{નાચવાનો} આનંદ માણે છે.

%%K t03-c2-13-1 p
13. --- છેડે, અમે \VUgras{અખાડો} \textsuperscript{(44)} જોઈએ છીએ,
જ્યાં બહાદુર સ્પેનિશ \VUgras{બળદલડવૈયો} ગુસ્સે થયેલા \VUgras{બળદ}
\textsuperscript{(45)} સાથે લડે છે ; પછી \VUgras{લપસણી}
\textsuperscript{(49)}, અને \VUgras{નિશાનબાજીની જગ્યા}
\textsuperscript{(46)}, જ્યાં લોકો \VUgras{બંદૂક} વાપરવાનો અને
\VUgras{નિશાન} પર તાકવાનો મહાવરો કરે છે.

%%K t03-c2-14-1 p
14. --- પાસે, આ રહ્યું \VUgras{મેળાનું નાટકઘર} \textsuperscript{(47)},
જ્યાં \VUgras{પ્રહસન} અને \VUgras{કરુણ નાટક} ભજવાય છે. તદ્દન પાસે,
\VUgras{રાક્ષસ} \textsuperscript{(48)} અને \VUgras{વામન}નો માંડવો છે.
પહેલાની \VUgras{ઊંચાઈ} બે મીટર પંદર સેન્ટિમીટર છે ; બીજો ભાગ્યે જ
બાળક જેટલો મોટો છે.

%%K t03-c2-15-1 p
15. --- છેલ્લે, આ રહ્યું મોટું \VUgras{પ્રાણીસંગ્રહાલય}
\textsuperscript{(50)} પોતાનાં \VUgras{જંગલી પ્રાણીઓ} સાથે, પોતાના
પ્રબળ \VUgras{નહોરવાળા} \VUgras{વાઘ} \textsuperscript{(51)} સાથે,
પોતાની ઘાટી \VUgras{કેશવાળીવાળા} \VUgras{સિંહ} \textsuperscript{(52)}
સાથે, પોતાનાં બે \VUgras{જિરાફ} \textsuperscript{(53)} અને પોતાના
\VUgras{હાથી} \textsuperscript{(54)} સાથે, જે પોતાની \VUgras{સૂંઢ} એટલી
કુશળતાથી વાપરે છે.

%%K t03-c2-16-1 p
16. --- આ પ્રાણીઓને કેળવનારો \VUgras{ખેલાડી} \textsuperscript{(55)}
માંડવાના બારણે છે.

\VUsaut{6.0mm}
\VUfilet{22mm}
